\subsection{Formalisme Mekanika Kuantum untuk Pemodelan Informasi}
Untuk membangun fondasi matematis yang diperlukan di dalam konsep Teori Permainan Kuantum dan Komputasi Kuantum, diberikan beberapa poin untuk memperdalam konsep yang linear dengan konsep utama \textbf{Fisika Kuantum} dan \textbf{Fisika Matematika}.

\subsubsection{Ruang Hilbert, Superposisi, dan Sistem Qubit-Qutrit}

Dalam upaya \textbf{prediksi pasar saham secara \textit{real-time}} \citep{bakshi_quantum_2025}, kita dihadapkan pada dinamika pasar yang sangat kompleks, seperti fluktuasi indeks S\&P500. Data pasar yang kaya ini menuntut representasi informasi yang padat. Namun, penggunaan sistem biner klasik atau bahkan sistem kuantum dua-level (\textbf{qubit}) sering kali mengalami keterbatasan. Keterbatasan utama muncul dari \textit{discretization error} saat merepresentasikan data kontinu, menyebabkan hilangnya informasi signifikan.

Fenomena ini bisa dianalogikan sebagaimana pada kasus \textbf{Quantum Penny-Flip Game} \citep{meyer_quantum_1999}, di mana pemain klasik terhambat oleh strategi biner (membalik atau tidak membalik koin), sedangkan pemain kuantum meraih kemenangan dengan memanfaatkan ruang keadaan yang lebih luas melalui superposisi. Dalam konteks \textit{financial forecasting}, untuk mengatasi keterbatasan representasi ini dan menangkap korelasi pasar yang kompleks dengan efisiensi komputasi yang lebih tinggi, kita membutuhkan kapasitas pemrosesan informasi yang lebih besar.

Solusi yang relevan adalah perluasan ruang keadaan dari qubit ke \textbf{qutrit} (sistem tiga-level). Qutrit menawarkan kepadatan informasi yang lebih tinggi ($\log_2 3 \approx 1.585$ bit per qutrit) and fleksibilitas topologi sirkuit yang lebih baik dibandingkan jaringan saraf berbasis qubit (QQBN). Transisi ke qutrit ini merupakan langkah fundamental yang kompleks untuk meningkatkan kapasitas pemrosesan informasi dalam aplikasi keuangan tingkat lanjut.

Secara formal, sistem kuantum dijelaskan dalam ruang vektor kompleks yang dilengkapi dengan \textit{inner product}, dikenal sebagai \textbf{Ruang Hilbert} ($\mathcal{H}$). Unit dasar informasi kuantum standar adalah \textbf{qubit}, dengan basis ortonormal \{$|0\rangle, |1\rangle$\}. State umum qubit dinyatakan sebagai:

\[ |\psi\rangle = \alpha|0\rangle + \beta|1\rangle \]

di mana $\alpha, \beta \in \mathbb{C}$ memenuhi $|\alpha|^2 + |\beta|^2 = 1$.

Selanjutnya, untuk sistem qutrit, yang beroperasi dalam ruang Hilbert tiga dimensi dengan basis komputasi \{$|0\rangle, |1\rangle, |2\rangle$\}, state superposisi qutrit dituliskan secara matematis sebagai:

\[ |\phi\rangle = c_0|0\rangle + c_1|1\rangle + c_2|2\rangle \]

dengan syarat normalisasi $\sum_{i=0}^2 |c_i|^2 = 1$.

\subsubsection{Formalisme Teori Permainan Klasik}

Sebelum membahas perluasan ke ranah kuantum, terdapat konsep penting yang perlu dibahas, yaitu mendefinisikan kerangka kerja teori permainan klasik secara formal. Sebuah permainan strategi murni dengan $n$-pemain didefinisikan oleh \textit{tuple} ruang strategi $\mathcal{S}$ and fungsi \textit{payoff} $\pi$:

\[ \mathcal{S} = (S_1, S_2, \dots, S_n) \]

\[ \pi = (\pi_1, \pi_2, \dots, \pi_n) \]

di mana $S_i = \{s_i^1, s_i^2, \dots \}$ adalah himpunan strategi murni pemain ke-$i$. Secara umum, permainan dalam bentuk normal (\textit{normal form}) dinyatakan sebagai \textit{ordered triple} $(N, (S_i)_{i \in N}, (u_i)_{i \in N})$ di mana $N$ adalah himpunan pemain dan $u_i$ adalah fungsi \textit{payoff} atau utilitas pemain $i$.

Solusi utama dalam teori permainan adalah \textbf{Nash Equilibrium}. Vektor strategi $s^* = (s_1^*, \dots, s_n^*)$ merupakan Nash Equilibrium jika tidak ada pemain yang dapat meningkatkan \textit{payoff}-nya dengan mengubah strategi secara sepihak:

\[ u_i(s^*) \geq u_i(s_i, s^*_{-i}) \quad \forall s_i \in S_i, \forall i \in N \]

Dalam notasi matriks, hal ini berarti:

\[ \pi_i(s_1^*, \dots, s_i^*, \dots, s_n^*) \geq \pi_i(s_1^*, \dots, s_i, \dots, s_n^*) \]

Konsep lain yang relevan adalah \textbf{Pareto Optimality}. Sebuah keadaan permainan $\mathcal{P}$ dikatakan mendominasi $\mathcal{P}'$ secara Pareto jika $\pi_i(\mathcal{P}) \geq \pi_i(\mathcal{P}')$ untuk semua pemain, with pertidaksamaan ketat berlaku setidaknya untuk satu pemain.

Dalam aplikasi ekonomi makro seperti \textbf{Barro-Gordon Game} \citep{lowe_linking_2024}, fungsi utilitas pembuat kebijakan ($U^{pol}$) sering dimodelkan sebagai fungsi dari inflasi ($\pi_t$) dan ekspektasi inflasi ($\pi_t^e$):

\[ U^{pol}_t = \theta b (\pi_t - \pi_t^e) - \frac{a \pi_t^2}{2} \]

Definisi-definisi klasik ini menjadi landasan untuk memahami bagaimana protokol kuantum memperluas ruang strategi dan keseimbangan yang dapat dicapai.

\subsubsection{Formalisme Matriks Densitas}

Dalam realitas pasar finansial, investor jarang memiliki informasi sempurna mengenai keadaan pasar. Pasar sering kali dipenuhi oleh \textit{noise} dan ketidakpastian yang membuat representasi sistem sebagai vektor keadaan murni (\textit{pure state}) $|\psi\rangle$ menjadi tidak realistis. Keadaan pasar lebih tepat digambarkan sebagai ansambel statistik dari berbagai kemungkinan keadaan, atau yang dikenal sebagai \textit{mixed state}.

Untuk mengakomodasi ketidakpastian intrinsik ini secara matematis, kita menggunakan \textbf{Matriks Densitas} ($\rho$). Formalisme ini memungkinkan kita untuk mendeskripsikan sistem di mana kita hanya mengetahui probabilitas $p_i$ bahwa sistem berada dalam keadaan kuantum tertentu $|\psi_i\rangle$. Definisi formalnya adalah:

\[ \rho = \sum_i p_i |\psi_i\rangle\langle\psi_i| \]

dengan syarat normalisasi $\text{Tr}(\rho) = 1$. Dalam konteks ini, elemen non-diagonal (koherensi) merepresentasikan superposisi kuantum yang dapat dimanfaatkan untuk keuntungan strategis, sementara elemen diagonal merepresentasikan probabilitas klasik. Evolusi sistem pasar yang dimodelkan ini kemudian mengikuti persamaan uniter $\rho_f = U \rho U^\dagger$, yang menjaga total probabilitas tetap satu.

\subsubsection{Protokol Kuantisasi Permainan}

Untuk mengkuantisasi permainan klasik, berbagai protokol telah dikembangkan untuk memetakan elemen-elemen permainan klasik ke dalam domain kuantum. Proses umum dari protokol tersebut melibatkan beberapa tahapan:

\begin{enumerate}
    \item \textbf{Sistem Kuantum}: Ruang strategi dan \textit{payoff} pemain klasik dipetakan ke dalam operator dan state dalam ruang Hilbert.
    \item \textbf{State Awal}: Permainan dimulai dengan state kuantum awal, seringkali merupakan state terjerat (\textit{entangled}) untuk mengeksploitasi fitur kuantum.
    \item \textbf{Strategi Pemain}: Strategi pemain direpresentasikan oleh operator uniter yang bekerja pada state kuantum.
    \item \textbf{Pengukuran dan Payoff}: Hasil permainan ditentukan oleh pengukuran state akhir, dan \textit{payoff} dihitung dari nilai ekspektasi operator \textit{payoff}.
\end{enumerate}

Berikut adalah dua protokol kuantisasi permainan yang umum:

\paragraph{Skema Eisert-Wilkens-Lewenstein (EWL)}

Protokol Eisert-Wilkens-Lewenstein (EWL) adalah salah satu kerangka kerja paling berpengaruh untuk mengkuantisasi permainan klasik. Protokol ini memperkenalkan operator \textit{entanglement} $J$ yang bekerja pada state awal klasik (misalnya, $|00\rangle$) untuk menciptakan korelasi kuantum sebelum pemain menerapkan strategi mereka \citep{eisert_quantum_1999, samadi_quantum_2017}.

Langkah-langkah utamanya adalah:

\begin{enumerate}
    \item \textbf{State Awal Terjerat}: State awal sistem disiapkan dalam bentuk terjerat menggunakan operator $J$:
    \[ |\psi_{in}\rangle = J |00\rangle \]
    Di mana $J$ sering mengambil bentuk $J = e^{i\gamma \sigma_x \otimes \sigma_x / 2}$, dengan $\gamma$ mengukur tingkat keterikatan.
    \item \textbf{Strategi Pemain}: Setiap pemain memilih strategi yang direpresentasikan oleh operator uniter lokal (misalnya, $U_A$ dan $U_B$).
    \item \textbf{State Akhir}: Setelah pemain menerapkan strategi, operator \textit{un-entanglement} $J^\dagger$ diterapkan untuk memisahkan state sebelum pengukuran:
    \[ |\psi_f\rangle = J^\dagger (U_A \otimes U_B) J |00\rangle \]
    \item \textbf{Payoff}: \textit{Payoff} dihitung dari hasil pengukuran state akhir ini. Protokol EWL memungkinkan tercapainya Nash Equilibrium baru yang seringkali Pareto-superior dibandingkan ekuilibrium klasik, terutama dalam permainan seperti Dilema Tahanan.
\end{enumerate}

\paragraph{Skema Marinatto-Weber (MW)}

Skema Marinatto-Weber (MW) menyediakan protokol standar lain untuk mengkuantisasi permainan klasik, yang relevan untuk situasi dengan \textit{mixed state} atau ketika state awal tidak murni \citep{marinatto_quantum_2000, samadi_quantum_2017}. Protokol ini sangat cocok untuk menganalisis \textit{Mixed Strategy Quantum Games}.

Langkah-langkah formalnya adalah:

\begin{enumerate}
    \item \textbf{State Awal}: Permainan dimulai dengan state awal $\rho_{in}$ pada ruang Hilbert gabungan $\mathcal{H}_A \otimes \mathcal{H}_B$, yang bisa berupa state terjerat (\textit{entangled}) atau state campuran.
    \item \textbf{Strategi}: Pemain A dan B memilih strategi yang direpresentasikan oleh operator uniter $U_A$ dan $U_B$. Aksi gabungan pada sistem adalah $U = U_A \otimes U_B$.
    \item \textbf{State Akhir}: State sistem setelah aksi pemain berevolusi menjadi:
    \[ \rho_f = U \rho_{in} U^\dagger = (U_A \otimes U_B) \rho_{in} (U_A^\dagger \otimes U_B^\dagger) \]
    \item \textbf{Payoff}: Payoff yang diterima pemain $i$ (misal, A) adalah nilai ekspektasi dari operator \textit{payoff} $P_A$:
    \[ \$_{A} = \text{Tr}(P_A \rho_f) \]
\end{enumerate}

Dalam konteks \textit{Quantum Barro-Gordon Game}, strategi klasik (inflasi tinggi/rendah) dipetakan ke operator Identitas ($I$) dan Pauli-X ($\sigma_x$) atau operator uniter lainnya, memungkinkan solusi \textit{Nash Equilibrium} yang \textit{time-consistent} melalui eksploitasi superposisi strategi.

\subsubsection{Teori Permainan Bayesian dan Quantum Discord}

Teori permainan klasik sering mengasumsikan informasi sempurna. Namun, dalam skenario nyata, pemain sering menghadapi ketidakpastian mengenai tipe lawan mereka. \textbf{Teori Permainan Bayesian} memodelkan situasi ini dengan menetapkan \textit{prior belief}.

Berdasarkan \citet{lowe_linking_2024}, payoff ekspektasi untuk pemain A dalam permainan Bayesian diberikan oleh:

\[ u_A = \sum_{\alpha, \beta} P_A(\alpha, \beta) \sum_{\sigma, \sigma'} P(\sigma \sigma' | \alpha \beta) u^{\alpha, \beta}_{\sigma, \sigma', A} \]

di mana $\alpha, \beta$ adalah tipe permainan/pemain, dan $\sigma, \sigma'$ adalah hasil pengukuran (aksi).

Keunggulan kuantum dalam skenario ini tidak hanya bergantung pada \textit{entanglement}, tetapi juga pada \textbf{Quantum Discord}. Quantum discord mengukur korelasi kuantum total dikurangi korelasi klasik, dan didefinisikan melalui perbedaan antara \textit{quantum mutual information} $I(\rho^{AB})$ dan korelasi klasik $J_A(\rho^{AB})$:

\[ D_A(\rho^{AB}) = I(\rho^{AB}) - J_A(\rho^{AB}) \]

\[ I(\rho^{AB}) = S(\rho^A) + S(\rho^B) - S(\rho^{AB}) \]

di mana $S(\rho)$ adalah entropi von Neumann. Quantum discord mampu menangkap korelasi kuantum non-lokal bahkan dalam \textit{mixed separable states}, yang krusial untuk aplikasi di mana entanglement murni sulit dipertahankan \citep{lowe_linking_2024}.

\subsubsection{Kognisi Kuantum dan Pelanggaran Prinsip Rasionalitas}

Teori keputusan klasik didasarkan pada aksioma rasionalitas, salah satunya adalah \textbf{Sure Thing Principle} dari Savage. Prinsip ini menyatakan bahwa jika seorang pengambil keputusan lebih memilih aksi $A$ daripada $B$ dalam keadaan $X$, dan juga lebih memilih $A$ daripada $B$ dalam keadaan bukan-$X$, maka ia harus memilih $A$ terlepas dari apakah $X$ terjadi atau tidak. Namun, eksperimen psikologi kognitif, seperti yang dilakukan oleh \citet{tversky_disjunction_1992} pada \textit{Prisoner's Dilemma}, menunjukkan pelanggaran prinsip ini yang dikenal sebagai \textbf{disjunction effect}.

Kognisi kuantum (\textit{Quantum Cognition}) memodelkan fenomena irasionalitas ini bukan sebagai kesalahan (error), melainkan sebagai manifestasi dari \textbf{Interferensi Kuantum}. Dalam formalisme ini, probabilitas total keputusan tidak mengikuti hukum aditivitas klasik (Hukum Total Probabilitas), melainkan mengandung suku interferensi:

\[ P(A) = | \psi_A + \psi_B |^2 = P(A) + P(B) + 2\sqrt{P(A)P(B)} \cos \theta \]

Di mana $\theta$ adalah fase relatif yang merepresentasikan bias kognitif atau ambiguitas. Jika $\cos \theta \neq 0$, probabilitas keputusan dapat menyimpang dari prediksi klasik, menjelaskan fenomena seperti \textit{panic selling} atau \textit{herd behavior} di pasar keuangan. 

Untuk memetakan struktur keputusan kompleks ini, digunakan \textbf{Quantum-like Influence Diagrams (QLID)}. QLID memperluas \textit{Bayesian Networks} klasik dengan mengganti probabilitas kondisional riil dengan amplitudo probabilitas kompleks, memungkinkan pemodelan superposisi keyakinan investor sebelum keputusan final (kolaps) diambil \citep{moreira_quantum-like_2020}.